% WACV 2027 Paper Template
% based on the ICCV 2025 template (https://media.eventhosts.cc/Conferences/ICCV2025/ICCV2025-Author-Kit-Feb.zip) with
% WACV-specific details (e.g., 2 tracks) from the WACV 2025 template (https://www.dropbox.com/scl/fi/su44zgdhrzik26p2xu37k/WACV-2025-Author-Kit-Template.zip?rlkey=5qcfimjhxnmx3wlyk7yhk8wg7&dl=0)

\documentclass[10pt,twocolumn,letterpaper]{article}

%%%%%%%%% PAPER TYPE  - PLEASE UPDATE FOR FINAL VERSION
% \usepackage[review,algorithms]{wacv}      % To produce the REVIEW version for the algorithms track
% \usepackage[review,applications]{wacv}      % To produce the REVIEW version for the applications track
\usepackage[review,datasets]{wacv}      % To produce the REVIEW version for the datasets track
% \usepackage{wacv}              % To produce the CAMERA-READY version
%\usepackage[pagenumbers]{wacv} % To force page numbers, e.g. for an arXiv version

% Import additional packages in the preamble file, before hyperref
\input{preamble}

% It is strongly recommended to use hyperref, especially for the review version.
% hyperref with option pagebackref eases the reviewers' job.
% Please disable hyperref *only* if you encounter grave issues, 
% e.g. with the file validation for the camera-ready version.
%
% If you comment hyperref and then uncomment it, you should delete *.aux before re-running LaTeX.
% (Or just hit 'q' on the first LaTeX run, let it finish, and you should be clear).
\definecolor{wacvblue}{rgb}{0.21,0.49,0.74}
\usepackage[pagebackref,breaklinks,colorlinks,allcolors=wacvblue]{hyperref}

%%%%%%%%% PAPER ID  - PLEASE UPDATE
\def\wacvPaperID{*****} % *** Enter the WACV Paper ID here
\def\confName{WACV}
\def\confYear{2027}


%%%%%%%%% TITLE - PLEASE UPDATE
\title{\LaTeX\ Author Guidelines for \confName~Proceedings}

%%%%%%%%% AUTHORS - PLEASE UPDATE
\author{First Author\\
Institution1\\
Institution1 address\\
{\tt\small firstauthor@i1.org}
% For a paper whose authors are all at the same institution,
% omit the following lines up until the closing ``}''.
% Additional authors and addresses can be added with ``\and'',
% just like the second author.
% To save space, use either the email address or home page, not both
\and
Second Author\\
Institution2\\
First line of institution2 address\\
{\tt\small secondauthor@i2.org}
}

\begin{document}
\maketitle
\input{sec/0_abstract}    
% =========================================================================
% PoleSight — paper outline / index for future writing.
% Scaffold only (WACV has no rigid section layout); expand each bullet into
% prose later. Section flow mirrors RoadSight (ACM MM '26) dataset paper.
% =========================================================================
% \section*{Outline / Index}
\section{Introduction}
\label{sec:intro}
\begin{itemize}
    \item \textbf{Why poles.} Pole-like objects (light, power, traffic, gantry,
        fence) are key roadside/utility infrastructure for mapping, asset
        inventory, and AV localization.
    \item \textbf{How sensed.} Captured by mobile laser scanners as LiDAR point
        clouds alongside camera images.
    \item \textbf{Range images work.} Projecting LiDAR to 2D range images lets
        mature 2D detectors handle pole detection efficiently.
    \item \textbf{Bottleneck.} Collecting + labeling real scans is slow and
        costly $\rightarrow$ synthetic data helps scale.
    \item \textbf{Simulation is costly.} Full LiDAR simulation needs both a 3D
        environment \emph{and} a sensor-simulation step.
    \item \textbf{Our idea.} Simple synthetic range-image generation
        (project $\rightarrow$ augment), no simulator $\rightarrow$ scalable datasets.
\end{itemize}

\paragraph{Contributions.}
\begin{itemize}
    \item \textbf{PoleSight} dataset: multi-modal range images, five pole
        classes, $\sim$2000 labeled frames.
    \item Scalable \textbf{LiDAR$\rightarrow$range-image} synthetic generation method.
    \item Baselines: range-image detection on public data + diverse LiDAR setups.
\end{itemize}

\section{Related Work}
\label{sec:related}
\begin{itemize}
    \item \textbf{Domain augmentation \& simulated LiDAR.} % TODO cite
    \item \textbf{Range-image learning.} RangeNet++, YOLO-family detectors. % TODO cite
    \item \textbf{Urban point clouds \& mobile LiDAR.} Pole extraction from MLS. % TODO cite
    \item \textbf{Pole / infrastructure datasets.} Position PoleSight vs.\ prior. % TODO cite
\end{itemize}

\section{Method}
\label{sec:method}

\subsection{Range-Image Generation}
\begin{itemize}
    \item Initial: per-frame range-image generation + labeling.
    \item Augmentation loop: projection $\rightarrow$ generation $\rightarrow$
        validation $\rightarrow$ repeat.
\end{itemize}

\subsection{Detection Model}
\begin{itemize}
    \item Architecture + training setup.
    \item Validation protocol.
\end{itemize}

\section{Dataset and Experiments}
\label{sec:dataset}

\subsection{Dataset}
\begin{itemize}
    \item $\sim$2000 labeled frames.
    \item Classes: fence, gantry-sign, light, power, traffic poles.
    \item Modalities: RGB, IR/intensity, LiDAR range.
\end{itemize}

\subsection{Results}
\begin{itemize}
    \item ML evaluation: detection metrics, per-class.
    \item Comparison vs.\ sequential point-cloud baselines.
\end{itemize}

\section{Discussion}
\label{sec:discussion}
\begin{itemize}
    \item \textbf{Strengths.} Works on public data; diverse LiDAR setups; beats
        sequential point-cloud processing.
    \item \textbf{Weaknesses.} Static objects only; computational cost.
    \item \textbf{Future.} Improved augmentation; mixed real + synthetic data.
\end{itemize}

\section{Conclusion}
\label{sec:conclusion}
% One paragraph: dataset + scalable synthetic range-image method, takeaway.
[00:42:43.656][Commander] BUILD command invoked.
[00:42:43.656][Build] The document of the active editor: file://%WS1%/main.tex
[00:42:43.656][Build] The languageId of the document: latex
[00:42:43.656][Root] Current workspace folders: ["file://%WS1%"]
[00:42:43.657][Root] Try finding root from magic comment.
[00:42:43.657][Root] Try finding root from active editor.
[00:42:43.657][Root] Found root file from active editor: %WS1%/main.tex
[00:42:43.657][Root] Keep using the same root file: %WS1%/main.tex
[00:42:43.657][Event] ROOT_FILE_SEARCHED
[00:42:43.658][Event] STRUCTURE_UPDATED
[00:42:43.658][Build] Building root file: %WS1%/main.tex
[00:42:43.658][Build][Recipe] Build root file %WS1%/main.tex
[00:42:43.663][Build][Recipe] Preparing to run recipe: latexmk.
[00:42:43.663][Build][Recipe] Prepared 1 tools.
[00:42:43.665][Build][Recipe] Cannot run `pdflatex` to determine if we are using MiKTeX. TypeError: Cannot read properties of undefined (reading 'toString')
[00:42:43.665]TypeError: Cannot read properties of undefined (reading 'toString')
	at isMikTeX (/Users/akash/.vscode/extensions/james-yu.latex-workshop-10.16.1/out/src/compile/recipe.js:393:80)
	at /Users/akash/.vscode/extensions/james-yu.latex-workshop-10.16.1/out/src/compile/recipe.js:369:63
	at Array.forEach (<anonymous>)
	at populateTools (/Users/akash/.vscode/extensions/james-yu.latex-workshop-10.16.1/out/src/compile/recipe.js:326:16)
	at createBuildTools (/Users/akash/.vscode/extensions/james-yu.latex-workshop-10.16.1/out/src/compile/recipe.js:170:5)
	at async build (/Users/akash/.vscode/extensions/james-yu.latex-workshop-10.16.1/out/src/compile/recipe.js:55:19)
	at async Object.build (/Users/akash/.vscode/extensions/james-yu.latex-workshop-10.16.1/out/src/compile/build.js:173:5)
	at async Object.build (/Users/akash/.vscode/extensions/james-yu.latex-workshop-10.16.1/out/src/core/commands.js:96:5)
	at async Pg._executeContributedCommand (file:///Applications/Visual%20Studio%20Code.app/Contents/Resources/app/out/vs/workbench/api/node/extensionHostProcess.js:534:48807)
[00:42:43.666][Build][Recipe] RootFile auxDir: %WS1% .
[00:42:43.676][Build] Recipe step 1 The command is latexmk:["-synctex=1","-interaction=nonstopmode","-file-line-error","-pdf","-outdir=%WS1%","-auxdir=%WS1%","%WS1%/main"].
[00:42:43.676][Build] env: undefined
[00:42:43.676][Build] root: %WS1%/main.tex
[00:42:43.676][Build] cwd: %WS1%
[00:42:43.677][Build] LaTeX build process spawned with PID undefined.
[00:42:43.678][Build] LaTeX fatal error on PID undefined. Error: spawn latexmk ENOENT
[00:42:43.678]Error: spawn latexmk ENOENT
	at ChildProcess._handle.onexit (node:internal/child_process:287:19)
	at onErrorNT (node:internal/child_process:508:16)
	at process.processTicksAndRejections (node:internal/process/task_queues:90:21)
[00:42:43.678][Build] Does the executable exist? $PATH: /Users/akash/.antigravity-ide/antigravity-ide/bin:/Users/akash/.antigravity/antigravity/bin:/Users/akash/.nvm/versions/node/v25.8.0/bin:/opt/homebrew/opt/ruby/bin:/opt/homebrew/lib/ruby/gems/4.0.0/bin:/opt/homebrew/opt/ruby/bin:/Applications/SnowSQL.app/Contents/MacOS:/Users/akash/.local/bin:/Users/akash/.pixi/bin:/usr/bin:/opt/homebrew/bin:/usr/local/opt/python/libexec/bin:/usr/local/bin:/System/Cryptexes/App/usr/bin:/usr/bin:/bin:/usr/sbin:/sbin:/var/run/com.apple.security.cryptexd/codex.system/bootstrap/usr/local/bin:/var/run/com.apple.security.cryptexd/codex.system/bootstrap/usr/bin:/var/run/com.apple.security.cryptexd/codex.system/bootstrap/usr/appleinternal/bin:/pkg/env/global/bin:/Library/Apple/usr/bin:/opt/homebrew/bin:/Users/akash/.lmstudio/bin, $Path: undefined, $SHELL: /bin/zsh
[00:42:43.678][Build] 

% =========================================================================
% PoleSight -- Related Work
% -------------------------------------------------------------------------
% Citation keys use AUTHOR-YEAR style. Every \cite{} key used below has a
% matching verification link + draft BibTeX entry in the REFERENCE MANIFEST
% at the bottom of this file (all commented out). Verify each link, then
% move the entries into main.bib (swap full venue names for the @String
% macros in main.bib if you prefer the abbreviated style).
% Only peer-reviewed conference / journal papers are cited.
% =========================================================================

\section{Related Work}
\label{sec:related}

We position PoleSight against four lines of work: synthetic and simulated
LiDAR for scaling supervision, range-image representations for learning on
LiDAR, mobile-laser-scanning analysis of urban scenes and pole extraction,
and existing pole or infrastructure datasets.

\paragraph{Domain augmentation and simulated LiDAR.}
Because annotating real LiDAR scans is slow and costly, a large body of work
generates synthetic point clouds to scale supervision. Driving simulators
such as CARLA~\cite{dosovitskiy2017carla} render virtual towns from which
sensor data can be sampled, and game-engine pipelines extract point-level
labels directly from virtual worlds~\cite{yue2018lidar}. To narrow the
appearance gap, LiDARsim~\cite{manivasagam2020lidarsim} couples physics-based
ray casting with a learned model of sensor noise to produce realistic scans
from real-world assets. Synthetic corpora such as
SynLiDAR~\cite{xiao2022synlidar} pair large simulated datasets with
point-cloud translation networks, while domain-adaptation frameworks like
ePointDA~\cite{zhao2021epointda} explicitly bridge the simulation-to-real
shift for segmentation. These approaches are effective but require a full 3D
environment together with a dedicated sensor-simulation or domain-adaptation
stage. In contrast, PoleSight scales data by augmenting range images
projected directly from real scans, avoiding both an explicit simulator and a
separate adaptation step.

\paragraph{Range-image learning.}
Projecting a rotating LiDAR's returns into a 2D range image lets mature image
CNNs and detectors operate on point clouds efficiently.
SqueezeSeg~\cite{wu2018squeezeseg} and RangeNet++~\cite{milioto2019rangenet}
perform real-time semantic segmentation in the range view, while
LaserNet~\cite{meyer2019lasernet} and RangeDet~\cite{fan2021rangedet} show
that range-view representations support competitive 3D object detection. This
line of work builds on single-stage 2D detectors such as
YOLO~\cite{redmon2016yolo}, whose speed makes them attractive for on-vehicle
deployment. Existing range-view methods, however, are developed and evaluated
almost exclusively on car, pedestrian, and cyclist classes from
autonomous-driving benchmarks. PoleSight instead targets thin, vertical
pole-like structures and provides a labeled range-image dataset that makes
such detectors directly applicable to roadside-infrastructure mapping.

\paragraph{Urban point clouds and mobile LiDAR.}
Mobile laser scanning (MLS) has long been used to inventory road
environments, and early pole-extraction methods rely on geometric
heuristics---voxel- or slice-based clustering and cylinder fitting to isolate
vertical structures~\cite{lehtomaki2010detection, pu2011recognizing,
cabo2014algorithm}. These rule-based pipelines are accurate in their target
scenes but are hand-tuned and do not scale to the variability of large,
learning-driven benchmarks. Large urban MLS and driving datasets---
Paris-Lille-3D~\cite{roynard2018parislille},
SemanticKITTI~\cite{behley2019semantickitti}, and
Toronto-3D~\cite{tan2020toronto}---provide dense point-wise annotations for
many urban categories and have driven progress in 3D deep learning. Yet poles
appear only as one incidental class among dozens, with limited per-type
granularity. PoleSight complements these resources with a pole-centric,
multi-class dataset designed for range-image detection rather than full-scene
3D segmentation.

\paragraph{Pole and infrastructure datasets.}
General autonomous-driving LiDAR benchmarks such as
KITTI~\cite{geiger2012kitti} and nuScenes~\cite{caesar2020nuscenes}, and urban
MLS datasets~\cite{roynard2018parislille, behley2019semantickitti,
tan2020toronto}, were built for vehicle-centric perception and treat
pole-like objects as a minor label, if at all. Dedicated infrastructure
datasets target narrower domains: TS40K~\cite{lavado2025ts40k}, for example,
captures electrical-transmission corridors from an aerial, top-down view of
rural terrain. To our knowledge, no public dataset offers dedicated,
multi-class roadside-pole annotations (light, power, traffic, gantry-sign,
and fence poles) in range-image form together with complementary RGB and
intensity channels. PoleSight fills this gap and pairs the dataset with a
scalable LiDAR-to-range-image synthetic-generation method and detection
baselines across public data and diverse LiDAR setups.


% =========================================================================
% REFERENCE MANIFEST -- verify each link, then add the entry to main.bib.
% (All venues are peer-reviewed conferences/journals. Lines flagged "VERIFY"
%  are best-effort author/title details to double-check against the link.)
% =========================================================================
%
% ---- Domain augmentation & simulated LiDAR ----
%
% dosovitskiy2017carla | CARLA: An Open Urban Driving Simulator | CoRL 2017 (PMLR v78)
%   link: https://proceedings.mlr.press/v78/dosovitskiy17a.html
% @inproceedings{dosovitskiy2017carla,
%   author    = {Dosovitskiy, Alexey and Ros, German and Codevilla, Felipe and Lopez, Antonio and Koltun, Vladlen},
%   title     = {{CARLA}: An Open Urban Driving Simulator},
%   booktitle = {Conference on Robot Learning (CoRL)},
%   pages     = {1--16}, year = {2017}}
%
% yue2018lidar | A LiDAR Point Cloud Generator: from a Virtual World to Autonomous Driving | ACM ICMR 2018
%   link: https://dl.acm.org/doi/10.1145/3206025.3206080
% @inproceedings{yue2018lidar,
%   author    = {Yue, Xiangyu and Wu, Bichen and Seshia, Sanjit A. and Keutzer, Kurt and Sangiovanni-Vincentelli, Alberto L.},
%   title     = {A {LiDAR} Point Cloud Generator: from a Virtual World to Autonomous Driving},
%   booktitle = {ACM International Conference on Multimedia Retrieval (ICMR)},
%   pages     = {458--464}, year = {2018}, doi = {10.1145/3206025.3206080}}
%
% manivasagam2020lidarsim | LiDARsim: Realistic LiDAR Simulation by Leveraging the Real World | CVPR 2020
%   link: https://openaccess.thecvf.com/content_CVPR_2020/html/Manivasagam_LiDARsim_Realistic_LiDAR_Simulation_by_Leveraging_the_Real_World_CVPR_2020_paper.html
% @inproceedings{manivasagam2020lidarsim,
%   author    = {Manivasagam, Sivabalan and Wang, Shenlong and Wong, Kelvin and Zeng, Wenyuan and Sazanovich, Mikita and Tan, Shuhan and Yang, Bin and Ma, Wei-Chiu and Urtasun, Raquel},
%   title     = {{LiDARsim}: Realistic {LiDAR} Simulation by Leveraging the Real World},
%   booktitle = {IEEE/CVF Conference on Computer Vision and Pattern Recognition (CVPR)},
%   pages     = {11167--11176}, year = {2020}}
%
% xiao2022synlidar | Transfer Learning from Synthetic to Real LiDAR Point Cloud for Semantic Segmentation | AAAI 2022
%   link: https://ojs.aaai.org/index.php/AAAI/article/view/20183
% @inproceedings{xiao2022synlidar,
%   author    = {Xiao, Aoran and Huang, Jiaxing and Guan, Dayan and Zhan, Fangneng and Lu, Shijian},
%   title     = {Transfer Learning from Synthetic to Real {LiDAR} Point Cloud for Semantic Segmentation},
%   booktitle = {AAAI Conference on Artificial Intelligence},
%   volume    = {36}, number = {3}, pages = {2795--2803}, year = {2022}}
%
% zhao2021epointda | ePointDA: An End-to-End Simulation-to-Real Domain Adaptation Framework for LiDAR Point Cloud Segmentation | AAAI 2021
%   link: https://ojs.aaai.org/index.php/AAAI/article/view/16464
% @inproceedings{zhao2021epointda,                                   % VERIFY full author list against link
%   author    = {Zhao, Sicheng and Wang, Yezhen and Li, Bo and Wu, Bichen and Gao, Yang and Xu, Pengfei and Darrell, Trevor and Keutzer, Kurt},
%   title     = {{ePointDA}: An End-to-End Simulation-to-Real Domain Adaptation Framework for {LiDAR} Point Cloud Segmentation},
%   booktitle = {AAAI Conference on Artificial Intelligence},
%   volume    = {35}, number = {4}, pages = {3500--3509}, year = {2021}}
%
% ---- Range-image learning ----
%
% wu2018squeezeseg | SqueezeSeg: Convolutional Neural Nets with Recurrent CRF for Real-Time Road-Object Segmentation from 3D LiDAR Point Cloud | ICRA 2018
%   link: https://doi.org/10.1109/ICRA.2018.8462926
% @inproceedings{wu2018squeezeseg,
%   author    = {Wu, Bichen and Wan, Alvin and Yue, Xiangyu and Keutzer, Kurt},
%   title     = {{SqueezeSeg}: Convolutional Neural Nets with Recurrent {CRF} for Real-Time Road-Object Segmentation from 3D {LiDAR} Point Cloud},
%   booktitle = {IEEE International Conference on Robotics and Automation (ICRA)},
%   pages     = {1887--1893}, year = {2018}, doi = {10.1109/ICRA.2018.8462926}}
%
% milioto2019rangenet | RangeNet++: Fast and Accurate LiDAR Semantic Segmentation | IROS 2019
%   link: https://doi.org/10.1109/IROS40897.2019.8967762
% @inproceedings{milioto2019rangenet,
%   author    = {Milioto, Andres and Vizzo, Ignacio and Behley, Jens and Stachniss, Cyrill},
%   title     = {{RangeNet++}: Fast and Accurate {LiDAR} Semantic Segmentation},
%   booktitle = {IEEE/RSJ International Conference on Intelligent Robots and Systems (IROS)},
%   pages     = {4213--4220}, year = {2019}, doi = {10.1109/IROS40897.2019.8967762}}
%
% meyer2019lasernet | LaserNet: An Efficient Probabilistic 3D Object Detector for Autonomous Driving | CVPR 2019
%   link: https://openaccess.thecvf.com/content_CVPR_2019/html/Meyer_LaserNet_An_Efficient_Probabilistic_3D_Object_Detector_for_Autonomous_Driving_CVPR_2019_paper.html
% @inproceedings{meyer2019lasernet,
%   author    = {Meyer, Gregory P. and Laddha, Ankit and Kee, Eric and Vallespi-Gonzalez, Carlos and Wellington, Carl K.},
%   title     = {{LaserNet}: An Efficient Probabilistic 3D Object Detector for Autonomous Driving},
%   booktitle = {IEEE/CVF Conference on Computer Vision and Pattern Recognition (CVPR)},
%   pages     = {12677--12686}, year = {2019}}
%
% fan2021rangedet | RangeDet: In Defense of Range View for LiDAR-based 3D Object Detection | ICCV 2021
%   link: https://www.computer.org/csdl/proceedings-article/iccv/2021/281200c898/1BmFeEDSC2I
% @inproceedings{fan2021rangedet,
%   author    = {Fan, Lue and Xiong, Xuan and Wang, Feng and Wang, Naiyan and Zhang, Zhaoxiang},
%   title     = {{RangeDet}: In Defense of Range View for {LiDAR}-based 3D Object Detection},
%   booktitle = {IEEE/CVF International Conference on Computer Vision (ICCV)},
%   pages     = {2918--2927}, year = {2021}}
%
% redmon2016yolo | You Only Look Once: Unified, Real-Time Object Detection | CVPR 2016
%   link: https://openaccess.thecvf.com/content_cvpr_2016/html/Redmon_You_Only_Look_CVPR_2016_paper.html
% @inproceedings{redmon2016yolo,
%   author    = {Redmon, Joseph and Divvala, Santosh and Girshick, Ross and Farhadi, Ali},
%   title     = {You Only Look Once: Unified, Real-Time Object Detection},
%   booktitle = {IEEE Conference on Computer Vision and Pattern Recognition (CVPR)},
%   pages     = {779--788}, year = {2016}}
%
% ---- Urban point clouds & mobile LiDAR (pole extraction) ----
%
% lehtomaki2010detection | Detection of Vertical Pole-Like Objects in a Road Environment Using Vehicle-Based Laser Scanning Data | Remote Sensing 2010
%   link: https://doi.org/10.3390/rs2030641
% @article{lehtomaki2010detection,
%   author  = {Lehtom{\"a}ki, Matti and Jaakkola, Anttoni and Hyypp{\"a}, Juha and Kukko, Antero and Kaartinen, Harri},
%   title   = {Detection of Vertical Pole-Like Objects in a Road Environment Using Vehicle-Based Laser Scanning Data},
%   journal = {Remote Sensing}, volume = {2}, number = {3}, pages = {641--664}, year = {2010}, doi = {10.3390/rs2030641}}
%
% pu2011recognizing | Recognizing basic structures from mobile laser scanning data for road inventory studies | ISPRS J. Photogramm. Remote Sens. 2011
%   link: https://www.sciencedirect.com/science/article/abs/pii/S0924271611000955
% @article{pu2011recognizing,
%   author  = {Pu, Shi and Rutzinger, Martin and Vosselman, George and Oude Elberink, Sander},
%   title   = {Recognizing basic structures from mobile laser scanning data for road inventory studies},
%   journal = {ISPRS Journal of Photogrammetry and Remote Sensing}, volume = {66}, number = {6}, pages = {S28--S39}, year = {2011}}
%
% cabo2014algorithm | An algorithm for automatic detection of pole-like street furniture objects from Mobile Laser Scanner point clouds | ISPRS J. Photogramm. Remote Sens. 2014
%   link: https://www.sciencedirect.com/science/article/pii/S092427161300230X
% @article{cabo2014algorithm,
%   author  = {Cabo, Carlos and Ord{\'o}{\~n}ez, Celestino and Garc{\'i}a-Cort{\'e}s, Silverio and Mart{\'i}nez, Javier},
%   title   = {An algorithm for automatic detection of pole-like street furniture objects from Mobile Laser Scanner point clouds},
%   journal = {ISPRS Journal of Photogrammetry and Remote Sensing}, volume = {87}, pages = {47--56}, year = {2014}, doi = {10.1016/j.isprsjprs.2013.10.008}}
%
% roynard2018parislille | Paris-Lille-3D: A large and high-quality ground-truth urban point cloud dataset for automatic segmentation and classification | IJRR 2018
%   link: https://doi.org/10.1177/0278364918767506
% @article{roynard2018parislille,
%   author  = {Roynard, Xavier and Deschaud, Jean-Emmanuel and Goulette, Fran{\c{c}}ois},
%   title   = {{Paris-Lille-3D}: A large and high-quality ground-truth urban point cloud dataset for automatic segmentation and classification},
%   journal = {International Journal of Robotics Research}, volume = {37}, number = {6}, pages = {545--557}, year = {2018}, doi = {10.1177/0278364918767506}}
%
% behley2019semantickitti | SemanticKITTI: A Dataset for Semantic Scene Understanding of LiDAR Sequences | ICCV 2019
%   link: https://openaccess.thecvf.com/content_ICCV_2019/html/Behley_SemanticKITTI_A_Dataset_for_Semantic_Scene_Understanding_of_LiDAR_Sequences_ICCV_2019_paper.html
% @inproceedings{behley2019semantickitti,
%   author    = {Behley, Jens and Garbade, Martin and Milioto, Andres and Quenzel, Jan and Behnke, Sven and Stachniss, Cyrill and Gall, J{\"u}rgen},
%   title     = {{SemanticKITTI}: A Dataset for Semantic Scene Understanding of {LiDAR} Sequences},
%   booktitle = {IEEE/CVF International Conference on Computer Vision (ICCV)},
%   pages     = {9297--9307}, year = {2019}}
%
% tan2020toronto | Toronto-3D: A Large-Scale Mobile LiDAR Dataset for Semantic Segmentation of Urban Roadways | CVPR Workshops 2020
%   link: https://openaccess.thecvf.com/content_CVPRW_2020/html/w11/Tan_Toronto-3D_A_Large-Scale_Mobile_LiDAR_Dataset_for_Semantic_Segmentation_of_CVPRW_2020_paper.html
% @inproceedings{tan2020toronto,
%   author    = {Tan, Weikai and Qin, Nannan and Ma, Lingfei and Li, Ying and Du, Jing and Cai, Guorong and Yang, Ke and Li, Jonathan},
%   title     = {{Toronto-3D}: A Large-Scale Mobile {LiDAR} Dataset for Semantic Segmentation of Urban Roadways},
%   booktitle = {IEEE/CVF Conference on Computer Vision and Pattern Recognition Workshops (CVPRW)},
%   pages     = {797--806}, year = {2020}}
%
% ---- Pole / infrastructure datasets (positioning) ----
%
% geiger2012kitti | Are we ready for autonomous driving? The KITTI Vision Benchmark Suite | CVPR 2012
%   link: https://doi.org/10.1109/CVPR.2012.6248074
% @inproceedings{geiger2012kitti,
%   author    = {Geiger, Andreas and Lenz, Philip and Urtasun, Raquel},
%   title     = {Are we ready for autonomous driving? The {KITTI} Vision Benchmark Suite},
%   booktitle = {IEEE Conference on Computer Vision and Pattern Recognition (CVPR)},
%   pages     = {3354--3361}, year = {2012}, doi = {10.1109/CVPR.2012.6248074}}
%
% caesar2020nuscenes | nuScenes: A Multimodal Dataset for Autonomous Driving | CVPR 2020
%   link: https://openaccess.thecvf.com/content_CVPR_2020/html/Caesar_nuScenes_A_Multimodal_Dataset_for_Autonomous_Driving_CVPR_2020_paper.html
% @inproceedings{caesar2020nuscenes,
%   author    = {Caesar, Holger and Bankiti, Varun and Lang, Alex H. and Vora, Sourabh and Liong, Venice Erin and Xu, Qiang and Krishnan, Anush and Pan, Yu and Baldan, Giancarlo and Beijbom, Oscar},
%   title     = {{nuScenes}: A Multimodal Dataset for Autonomous Driving},
%   booktitle = {IEEE/CVF Conference on Computer Vision and Pattern Recognition (CVPR)},
%   pages     = {11621--11631}, year = {2020}}
%
% lavado2025ts40k | TS40K: a 3D Point Cloud Dataset of Rural Terrain and Electrical Transmission System | WACV 2025
%   link: https://openaccess.thecvf.com/content/WACV2025/papers/Lavado_Learning_under_Noisy_Labels_Spurious_Points_and_Diverse_Structures_TS40K_WACV_2025_paper.pdf
% @inproceedings{lavado2025ts40k,                                    % VERIFY exact title/authors against link
%   author    = {Lavado, Diogo and Soares, Cl{\'a}udia and Micheletti, Alessandra and others},
%   title     = {{TS40K}: a 3D Point Cloud Dataset of Rural Terrain and Electrical Transmission System},
%   booktitle = {IEEE/CVF Winter Conference on Applications of Computer Vision (WACV)},
%   year      = {2025}}
%
% =========================================================================

\input{sec/1_intro}
\input{sec/2_formatting}
\input{sec/3_finalcopy}
{
    \small
    \bibliographystyle{ieeenat_fullname}
    \bibliography{main}
}


\end{document}
